\section{同余方程}

\subsection{作业}

\begin{problem}
	习题 1
	\begin{solution}
		\textbf{中国剩余定理}：设 $m_1, m_2, \dots, m_k$ 是两两互素的正整数，对于同余方程组
		$$
		x \equiv a_i \pmod{m_i} \quad,\ i = 1, 2, \dots, k
		$$
		在模 $M = \prod_{i = 1}^{k} m_i$ 意义下有唯一解。

		\textbf{证明}：分为两部分证明。
		\begin{itemize}
			\item \textbf{存在性}：设 $M_i = \frac{M}{m_i}$，可得 $M_i$ 与 $m_i$ 互素。因此存在
			$$
			y_i = M_i^{-1} \bmod m_i
			$$
			那么构造出解
			$$
			\hat{x} = \sum_{i = 1}^{k} a_i M_i y_i
			$$
			因为
			$$
			a_i M_i y_i \bmod m_j = \begin{dcases}
				0,\ i \neq j \\
				a_i,\ i = j
			\end{dcases}
			$$
			因此 $\hat{x}$ 确实是一个解。

			\item \textbf{唯一性}：若在模 $M$ 意义下有两个解 $x_1, x_2$（$0 \le x_1 < x_2 < M$），那么令 $d = x_2 - x_1$，则 $d$ 满足：
			$$
			d \equiv 0 \pmod{a_i} \quad,\ i = 1, 2, \dots, k
			$$
			因此可知 $\lcm(a_1, a_2, \dots, a_k) \mid d$，而 $d > 0$，则 $d \ge M$。这与假设矛盾。
		\end{itemize}

		\textbf{应用案例}：在 RSA 加密中利用中国剩余定理进行密钥验证。
	\end{solution}
\end{problem}

\begin{problem}
	习题 2, 3
	\begin{solution}
		\ 
		\begin{center}
			\large \textbf{《How to Share a Secret》读书笔记}
		\end{center}

		本论文由 Adi Shamir 撰写，旨在解决密码学系统中的密钥管理问题 。作者提出了一种 $(k, n)$ 阈值方案，将秘密数据 $D$ 分割成 $n$ 个部分$D_{1},...,D_{n}$。该方案确保只要掌握其中任意 $k$ 个或更多的部分，就能轻松计算出原始数据 $D$ ；但如果只掌握 $k-1$ 个或更少的部分，则无法获取关于 $D$ 的任何信息，所有可能的数值都具有同等概率。这种机制在提高系统可靠性的同时，避免了因单点被破译而导致的信息泄露风险。

		该方案基于多项式插值原理。在二维平面上，给定 $k$ 个横坐标互不相同的点，能够唯一确定一个次数为 $k-1$ 的多项式 $q(x)$ 。

		为了确保绝对安全，计算是在模一个素数 $p$ 的有限域上进行的，其中 $p$ 必须大于数据 $D$ 和分块数 $n$。具体构造如下：构建一个 $k-1$ 次多项式：
		$$
		q(x)=a_{0}+a_{1}x+...+a_{k-1}x^{k-1}
		$$
		其中将秘密数据 $D$ 作为常数项，即 $a_{0}=D$。其余系数从 $[0, p)$ 的均匀分布中随机选取。随后，计算出 $n$ 个分块的值，所有计算均取模 $p$：
		$$
		D_{1}=q(1),...,D_{i}=q(i),...,D_{n}=q(n)
		$$
		持有任意 $k$ 个分块的人，可以通过插值法求出 $q(x)$ 的各项系数，进而通过计算 $q(0)$ 恢复出秘密 $D$。而如果只有 $k-1$ 个分块，对于任意一个可能的候选值 $D'$，都能构造出唯一一个经过这 $k-1$ 个已知点且满足 $q'(0)=D'$ 的 $k-1$ 次多项式，因此攻击者完全无法推断出真实的 $D$ 值。

		此方案具有极高的实用性。首先，每个分块的大小不会超过原始数据的大小。其次，在不改变阈值 $k$ 的情况下，可以动态地添加或删除分块，无需影响其他现有的分块。此外，只需更换一个拥有相同常数项的新多项式，即可在不改变原始秘密 $D$ 的情况下更新所有分块，防止被暴露的分块随着时间积累而导致泄密。最后，通过为不同级别的人员分配不同数量的分块，该方案能够实现层级化的权限管理。
	\end{solution}
\end{problem}

\begin{problem}
	习题 4
	\begin{solution}
		\textbf{收获}：我从本课程中获得了很多线性代数在各个领域的具体应用，增加了我对不同领域的了解。比如系统控制论的几节课我学习到了控制论的一些基础知识，虽然不见得能够精通，但是增长了见闻。

		\textbf{建议}：部分课程（比如系统控制理论）中内容很有难度，但是实际上讲起来则跳过了很多东西，浮光掠影。希望课程能够尽量讲深讲透，即使内容少一些简单一些。
	\end{solution}
\end{problem}